\cvsection{Expériences professionnelles}

\begin{cventries}
  \cventry
    {Stagiaire ingénieur de recherche en IA}
    {INRIA -- Équipe Mnémosyne}
    {Bordeaux, France}
    {Mars 2026 -- Septembre 2026}
    {
      \begin{cvitems}
        \item {Conception et développement en autonomie, avec le suivi de mes encadrants, d’une application et d’un agent IA multi-étapes pour la recherche d’articles scientifiques.}
        \item {Construction de l’agent en Python sans framework d’orchestration afin d’en comprendre les mécanismes, avec API OpenAI, appels d’outils et gestion de l’état et de la mémoire.}
        \item {Conception de prompts adaptés au cas d’usage scientifique et à l’utilisation des outils, avec gestion du contexte et sélection des messages par l’utilisateur.}
        \item {Définition d’un protocole expérimental pour comparer Qwen, Gemma et GLM selon la qualité des réponses, la latence, le TTFT et la taille du contexte.}
        \item {Comparaison de plusieurs quantifications pour analyser les compromis entre qualité et performances selon les environnements d’exécution.}
        \item {Déploiement et serving de LLM avec vLLM sur une infrastructure HPC équipée de GPU NVIDIA, avec utilisation de Slurm pour les expérimentations.}
        \item {Analyse des résultats et rédaction en cours d’un article scientifique sur la comparaison des modèles de langage pour notre contexte.}
        \item {Mise en œuvre d’un pipeline RAG avec chunking, embeddings et recherche sémantique dans ChromaDB, associé à MongoDB.}
        \item {Développement d’API REST Python/FastAPI documentées avec OpenAPI et d’une interface React/Next.js pour l’utilisation de l’agent.}
        \item {Conteneurisation de composants avec Docker, versionnement sous Git et mise en place d’une CI/CD sur GitHub.}
      \end{cvitems}
    }

  \cventry
    {Stagiaire de recherche en apprentissage par renforcement}
    {Université Karunya}
    {Coimbatore, Inde}
    {Juin 2025 -- Sept. 2025}
    {
      \begin{cvitems}
        \item {Développement d’un modèle d’apprentissage par renforcement et de son environnement d’entraînement pour la navigation autonome d’un robot hospitalier.}
        \item {Déploiement et optimisation du logiciel et du modèle sur Raspberry Pi pour des tests en conditions réelles.}
      \end{cvitems}
    }
\end{cventries}
